\documentclass[11pt,a4paper,fontset=windows]{ctexart}
% ===== 公共导言区：高等代数（一）期中模拟卷 =====
\usepackage{amsmath,amssymb}
\usepackage{unicode-math}
\setmathfont{TeX Gyre Termes Math}
\usepackage[margin=2.05cm,top=2.35cm,bottom=2.0cm,headsep=11pt,includefoot]{geometry}
\usepackage{fancyhdr}
\usepackage{lastpage}
\usepackage{enumitem}
\usepackage{array}
\usepackage{bm}
\usepackage{xcolor}
\usepackage{zhnumber}
\usepackage{needspace}

\definecolor{hdgray}{gray}{0.32}
\definecolor{lngray}{gray}{0.5}
\definecolor{notebg}{gray}{0.955}

\pagestyle{fancy}
\fancyhf{}
\renewcommand{\headrulewidth}{0.4pt}
\renewcommand{\footrulewidth}{0pt}
\fancyhead[L]{\small\color{hdgray}\songti 高等代数（一）\ 期中模拟试卷（\thepp）}
\fancyhead[R]{\small\color{hdgray}\songti 考试时间：120 分钟\quad 满分：100 分}
\fancyfoot[C]{\small\color{lngray}\songti 第 \thepage\ 页\quad 共 \pageref{LastPage}\ 页}

\setlength{\parindent}{0pt}
\setlength{\parskip}{2pt}
\linespread{1.16}

% 题目列表 1—8 连续编号
\newlist{qs}{enumerate}{1}
\setlist[qs]{label=\textbf{\arabic*.},ref=\arabic*,leftmargin=2.55em,
  labelsep=0.55em,itemsep=23pt,topsep=6pt,parsep=0pt}
\newlist{note}{enumerate}{1}
\setlist[note]{label=(\arabic*),leftmargin=1.75em,itemsep=1.5pt,topsep=2.5pt,parsep=0pt}
\newlist{sub}{enumerate}{1}
\setlist[sub]{label=(\arabic*),leftmargin=1.9em,itemsep=3pt,topsep=3pt,parsep=0pt}

\newcommand{\secA}[1]{\Needspace{9\baselineskip}\par\vspace{10pt}{\large\bfseries\songti #1}\par\vspace{3pt}\nopagebreak}
\newcommand{\blank}[1][2.6cm]{\underline{\hspace{#1}}}
\newcounter{pp}
\renewcommand{\thepp}{\chinese{pp}}

\setcounter{pp}{2}
\begin{document}

\begin{center}
{\LARGE\bfseries\songti 高等代数（一）期中模拟试卷（二）}\\[7pt]
{\normalsize\songti 适用范围：第一章\ 多项式（\S1\,--\,\S9）；第二章\ 行列式（\S2.1\,--\,\S2.7）}\\[3pt]
{\small\color{hdgray}\songti 多元多项式、对称多项式不在考查范围}
\end{center}
\vspace{-4pt}

\begin{center}
\begin{tabular}{@{}p{0.335\textwidth}p{0.335\textwidth}p{0.26\textwidth}@{}}
\songti 姓名：\blank[2.5cm] & \songti 学号：\blank[2.5cm] & \songti 得分：\blank[1.9cm]
\end{tabular}
\end{center}

\vspace{2pt}
\hrule height 0.5pt
\begin{center}
\setlength{\fboxsep}{6pt}%
\colorbox{notebg}{\parbox{0.94\textwidth}{\small\songti
\textbf{考生须知：}
\begin{note}
\item 本试卷共 8 题，满分 100 分，考试时间 120 分钟。
\item 计算题须写出必要的计算过程，证明题须给出完整的推理步骤；只写答案不给分。
\item 所有答案一律写在答题纸上，写在试卷上无效。
\end{note}}}
\end{center}
\hrule height 0.5pt
\vspace{4pt}

\secA{一、计算题（每小题 10 分，共 40 分）}
\begin{qs}
\item 设 $n\ge 3$，$x_1,x_2,\ldots,x_n$ 为任意数。计算 $n$ 阶行列式
\[
D_n=\begin{vmatrix}
1&1&\cdots&1\\
x_1&x_2&\cdots&x_n\\
\vdots&\vdots&&\vdots\\
x_1^{\,n-2}&x_2^{\,n-2}&\cdots&x_n^{\,n-2}\\
x_1^{\,n}&x_2^{\,n}&\cdots&x_n^{\,n}
\end{vmatrix},
\]
并将结果写成 $n$ 阶范德蒙德行列式与 $\displaystyle\sum_{i=1}^{n}x_i$ 的乘积。

\item 设 $a_1,\ldots,a_n$ 与 $b_1,\ldots,b_n$ 为任意数，计算 $n$ 阶行列式
\[
D_n=\begin{vmatrix}
a_1+b_1&a_1+b_2&\cdots&a_1+b_n\\
a_2+b_1&a_2+b_2&\cdots&a_2+b_n\\
\vdots&\vdots&&\vdots\\
a_n+b_1&a_n+b_2&\cdots&a_n+b_n
\end{vmatrix},
\]
并就 $n=2$ 与 $n\ge 3$ 两种情形分别写出结果。

\item 设
\[
f(x)=x^{5}-x^{4}-2x^{3}+2x^{2}+x-1 .
\]
求 $(f(x),\,f'(x))$，并写出 $f(x)$ 的所有重因式及其重数。

\item 设 $f(x)=x^{3}+ax+b\in\mathbb{R}[x]$，其中 $a,b$ 为实参数。
\begin{sub}
\item 求使 $f(x)$ 有重根的全部实数对 $(a,b)$；
\item 对满足条件的 $(a,b)$，写出 $f(x)$ 的重根。
\end{sub}
\end{qs}

\secA{二、证明题（每小题 15 分，共 60 分）}
\begin{qs}[resume*]
\item 设 $m,n,k$ 为非负整数。证明：
\[
x^{2}+x+1\ \Big|\ x^{3m}+x^{3n+1}+x^{3k+2}.
\]

\item 设 $f(x)\in\mathbb{R}[x]$，$\deg f=n\ge 1$。证明：
\begin{sub}
\item 若 $\alpha\in\mathbb{C}$ 是 $f(x)$ 的 $k$ 重根，则共轭数 $\bar\alpha$ 也是 $f(x)$ 的 $k$ 重根；
\item 若 $n$ 为奇数，则 $f(x)$ 在 $\mathbb{R}$ 中至少有一个根；
\item 若 $f(x)$ 在 $\mathbb{R}[x]$ 中不可约且 $\deg f\ge 2$，则 $\deg f=2$。
\end{sub}

\item 设 $p$ 为素数，$\Phi_p(x)=x^{p-1}+x^{p-2}+\cdots+x+1$。证明：$\Phi_p(x)$ 在 $\mathbb{Q}[x]$ 上不可约。

\item 设 $a_1,a_2,\ldots,a_n$ 是数域 $P$ 中 $n$ 个互不相同的数，
\[
V_n=\begin{vmatrix}
1&1&\cdots&1\\
a_1&a_2&\cdots&a_n\\
a_1^{2}&a_2^{2}&\cdots&a_n^{2}\\
\vdots&\vdots&&\vdots\\
a_1^{\,n-1}&a_2^{\,n-1}&\cdots&a_n^{\,n-1}
\end{vmatrix}.
\]
\begin{sub}
\item 证明 $V_n=\displaystyle\prod_{1\le i<j\le n}(a_j-a_i)\neq 0$；
\item 设 $b_1,\ldots,b_n\in P$。证明：存在唯一的多项式 $f(x)\in P[x]$，$\deg f\le n-1$，使得
$f(a_i)=b_i\ (i=1,2,\ldots,n)$。
\end{sub}
\end{qs}

\end{document}
